% =====================================================================
%  Page de titre — conforme aux directives EPFL / Section IC
%
%  À COMPLÉTER AVANT SOUMISSION : les champs entre crochets ci-dessous
%  sont des placeholders. Les directives de la section exigent :
%    - le titre                                          [rempli]
%    - les coordonnées de l'étudiant (nom, prénom, adresse)  [adresse à remplir]
%    - le nom du laboratoire EPFL / université / entreprise  [à remplir]
%    - le nom du superviseur académique EPFL responsable     [à remplir]
% =====================================================================
\begin{titlepage}
\thispagestyle{empty}
\centering

\vspace*{0.5cm}

{\Large\scshape École Polytechnique Fédérale de Lausanne\par}
\vspace{0.3cm}
{\large Faculté Informatique et Communications\par}
\vspace{0.15cm}
{\large Section Data Science\par}

\vspace{2.2cm}

{\large\bfseries Thèse de master\par}
\vspace{0.4cm}

\vspace{1.6cm}

\rule{\textwidth}{0.8pt}
\vspace{0.7cm}

{\huge\bfseries DOM-Graph RAG Visuel\par}
\vspace{0.5cm}
{\LARGE Tuilage guidé par la structure du document\\[0.25cm]
pour le RAG visuel\par}

\vspace{0.7cm}
\rule{\textwidth}{0.8pt}

\vspace{2.2cm}

{\Large\bfseries Alexis Firome\par}
\vspace{0.35cm}
{\normalsize
\begin{tabular}{r@{\hspace{0.6em}}l}
Courriel : & \href{mailto:alexis.firome@epfl.ch}{alexis.firome@epfl.ch}\\[0.2cm]
Adresse : & {27 ter boulevard Diderot, 75012 Paris}\\
\end{tabular}\par}

\vfill

{\normalsize
\begin{tabular}{@{}>{\raggedright\arraybackslash}p{6.0cm}@{\hspace{0.4cm}}>{\raggedright\arraybackslash}p{7.6cm}@{}}
\textbf{Superviseur académique EPFL} & {Pr. Amir Zamir}\\[0.35cm]
\textbf{Laboratoire d'accueil} & {Devoteam DRI}\\[0.35cm]
\textbf{Encadrement} & {Dr. Valentin Noël}\\[0.35cm]
\textbf{Date de soumission} & {09-09-2026}\\
\end{tabular}\par}

\vspace{1.4cm}

%{\small\itshape
%Ce travail a été réalisé conformément aux directives de la Section Informatique et Communications relatives au projet de master, et dans le respect des règles d'éthique de l'EPFL en matière de plagiat.\par}

\vspace{0.8cm}

\end{titlepage}

% ---------------------------------------------------------------------
\cleardoublepage
\thispagestyle{empty}
\vspace*{2cm}
\begin{center}
{\Large\bfseries Déclaration}
\end{center}
\vspace{1cm}

\noindent
Je déclare que ce rapport présente le travail que j'ai personnellement réalisé
dans le cadre de mon projet de master. Les contributions de tiers, les
résultats issus de la littérature et les outils logiciels réutilisés sont
explicitement identifiés et référencés. Le code source du système décrit dans
ce rapport est disponible en accès ouvert.

\vspace{1cm}

\noindent
Une version antérieure et condensée des travaux décrits aux
\cref{chap-methodologie,chap-resultats} a été soumise sous forme d'article
court à une conférence en traitement automatique des langues~; le présent
rapport en constitue la version longue, complétée du contexte scientifique
(\cref{chap-etat_de_lart}), du détail d'implémentation et de l'analyse
critique des résultats.

\vspace{2.5cm}

\noindent
\begin{tabular}{@{}p{7cm}p{7cm}@{}}
{Paris}, le {09-09-2026} & \hfill Alexis Firome\\[1.6cm]
& \hrulefill\\
& \hfill\small signature\\
\end{tabular}

\vspace{\fill}

% ---------------------------------------------------------------------
\cleardoublepage
\thispagestyle{empty}
\vspace*{3cm}
\begin{center}
{\Large\bfseries Remerciements}
\end{center}
\vspace{1cm}

\noindent
\textcolor{red}{[Section à personnaliser avant soumission.]}

\vspace{0.5cm}

\noindent
Je remercie \textcolor{red}{[nom du superviseur]} pour l'encadrement de ce
projet, ainsi que \textcolor{red}{[laboratoire / personnes]} pour les
discussions techniques et l'accès aux ressources de calcul qui ont rendu
possible la partie expérimentale de ce travail. Les erreurs et les choix
discutables qui subsistent dans ce rapport sont les miens.

\vspace{\fill}
